\documentclass[11pt,a4paper]{article}

\usepackage[T1]{fontenc}
\usepackage[utf8]{inputenc}
\usepackage[polish]{babel}
\usepackage[margin=2.15cm,top=2.55cm,bottom=2.25cm,headheight=22pt,headsep=0.65cm]{geometry}
\usepackage{graphicx}
\usepackage[table]{xcolor}
\usepackage{amsmath,amssymb}
\usepackage{array,booktabs,longtable}
\usepackage{enumitem}
\usepackage{fancyhdr}
\usepackage{titlesec}
\usepackage{microtype}
\usepackage{hyperref}

\definecolor{PWForest}{HTML}{0D2A00}
\definecolor{PWLeaf}{HTML}{778C12}
\definecolor{PWField}{HTML}{385A15}
\definecolor{PWDeepGreen}{HTML}{103101}
\definecolor{PWOrange}{HTML}{EE8106}
\definecolor{PWIvory}{HTML}{F7F6EE}

\setlength{\parindent}{0pt}
\setlength{\parskip}{0.55\baselineskip}
\setlength{\emergencystretch}{3em}
\setlist{topsep=0.25\baselineskip,itemsep=0.2\baselineskip,leftmargin=1.6em}
\setlist[itemize]{label=\textcolor{PWOrange}{\textbullet}}
\setlist[enumerate]{label=\colorbox{PWForest}{\color{white}\strut\arabic*},leftmargin=2em}
\renewcommand{\arraystretch}{1.3}
\setlength{\LTpre}{0.55\baselineskip}
\setlength{\LTpost}{0.55\baselineskip}
\newcommand{\code}[1]{\textcolor{PWField}{\texttt{\detokenize{#1}}}}

\title{Pole Widzenia\\[0.4em]\large Dane, które dają plony.}
\author{Zespół Jakobiany\\HackNation 2026 -- wyzwanie KOWR}
\date{}

\hypersetup{
  colorlinks=true,
  linkcolor=PWField,
  urlcolor=PWOrange,
  pdftitle={Pole Widzenia -- Dane, które dają plony},
  pdfauthor={Zespół Jakobiany}
}

\titleformat{\section}
  {\sffamily\Large\bfseries\color{PWForest}}
  {\colorbox{PWForest}{\color{white}\strut\thesection}}
  {0.65em}
  {}
  [\vspace{0.35ex}{\color{PWOrange}\titlerule}]
\titlespacing*{\section}{0pt}{2.0ex plus 0.7ex minus 0.2ex}{1.0ex}
\titleformat{\subsection}
  {\sffamily\large\bfseries\color{PWField}}
  {\thesubsection}
  {0.65em}
  {}
\titlespacing*{\subsection}{0pt}{1.5ex plus 0.5ex minus 0.2ex}{0.5ex}

\pagestyle{fancy}
\fancyhf{}
\fancyhead[L]{\sffamily\footnotesize\bfseries\color{PWForest} POLE WIDZENIA}
\fancyhead[R]{\sffamily\footnotesize\color{PWField} NOTA PROJEKTOWA}
\fancyfoot[L]{\sffamily\footnotesize\color{PWField} Zespół Jakobiany\textcolor{PWOrange}{\enspace\textbullet\enspace} HackNation 2026}
\fancyfoot[R]{\sffamily\footnotesize\bfseries\color{PWForest}\thepage}
\renewcommand{\headrulewidth}{0.5pt}
\renewcommand{\headrule}{\hbox to\headwidth{\color{PWLeaf}\leaders\hrule height \headrulewidth\hfill}}
\renewcommand{\footrulewidth}{0pt}

\begin{document}
\pagecolor{PWIvory}
\hypersetup{pageanchor=false}
\begin{titlepage}
\thispagestyle{empty}
\vspace*{0.8cm}
{\color{PWOrange}\rule{1.5cm}{3pt}}\par
\vspace{0.45cm}
{\sffamily\footnotesize\bfseries\color{PWField} HACKNATION 2026\enspace /\enspace WYZWANIE KOWR\par}
\vspace{1.35cm}
\begin{center}
  \includegraphics[width=0.66\textwidth]{pole_widzenia_z_napisem_przezroczyste.png}
\end{center}
\vspace{1.1cm}
{\sffamily\bfseries\color{PWForest}\fontsize{25}{31}\selectfont
Dane, które dają plony.\\ \par}
\vspace{0.8cm}
\noindent\color{PWLeaf}\rule{0.52\textwidth}{2pt}\hfill
\color{PWOrange}\rule{0.16\textwidth}{2pt}\hfill
\color{PWField}\rule{0.22\textwidth}{2pt}\par
\vfill
\noindent\colorbox{PWForest}{%
  \parbox{\dimexpr\textwidth-2\fboxsep\relax}{%
    \vspace{0.35cm}
    {\sffamily\small\bfseries\color{white} ZESPÓŁ JAKOBIANY\par}
    \vspace{0.12cm}
    {\sffamily\small\color{PWIvory} HackNation 2026 -- wyzwanie KOWR\par}
    \vspace{0.35cm}
  }%
}\par
\vspace*{0.7cm}
\end{titlepage}
\hypersetup{pageanchor=true}
\nopagecolor
\clearpage

\section{Opis projektu}

Pole Widzenia to aplikacja webowa wspierająca rolnika w ocenie sytuacji wodnej wybranej działki i planowaniu działań retencyjnych. Łączy informacje pogodowe, glebowe, satelitarne i terenowe, aby oszacować zapas wody w strefie korzeniowej, uzasadnić rekomendacje i wskazać, co warto zweryfikować w terenie.

\subsection{Problem i cel}

Dane przydatne do gospodarowania wodą są rozproszone między różnymi usługami i zbiorami. Pole Widzenia agreguje je dla wybranej działki, daty i uprawy oraz przedstawia wynik w raporcie z informacją o pochodzeniu, dostępności i ograniczeniach danych. Aplikacja ma wspierać decyzje rolnika, a nie zastępować pomiary polowe ani samodzielną ocenę warunków gospodarstwa.

\subsection{Funkcjonalności}

\begin{itemize}
  \item Wyszukiwanie działki EGiB po identyfikatorze, nazwie lub obrębie, wskazanie punktu na mapie, narysowanie granicy albo import GeoJSON (\code{Polygon} lub \code{MultiPolygon} w WGS 84).
  \item Analiza daty z ostatnich dziesięciu lat dla sześciu profili: pszenicy ozimej, kukurydzy na ziarno, ziemniaka, rzepaku ozimego, buraka cukrowego i użytków zielonych. Dostępny jest także profil ogólny.
  \item Podanie dodatkowego kontekstu: daty siewu, fazy rozwojowej, znanego nawodnienia i obserwacji wilgotności.
  \item Raport z bilansem pogodowym, porównaniami sezonowymi, szacowanym zapasem wody w strefie korzeniowej oraz kontekstem satelitarnym i terenowym. Rekomendacje zawierają uzasadnienie i mogą obejmować kontrolę warunków w terenie.
  \item Paszport pola z danymi o działce, glebie, pogodzie, roślinności i terenie, wraz z metadanymi pochodzenia i statusami dostępności danych.
  \item Prezentacja granicy działki, wykresów i tabel oraz eksport raportu do PDF lub JSON. Paszport i granicę można również wyeksportować odpowiednio jako JSON i GeoJSON.
  \item Instrukcje prostych eksperymentów terenowych: oceny wilgotności gleby, wsiąkania, struktury gleby metodą VESS i rozpadu bryłek w wodzie.
\end{itemize}

\section{Sposób analizy}

\begin{enumerate}
  \item Użytkownik wskazuje geometrię działki, datę i profil uprawy; może dodać informacje znane z gospodarstwa.
  \item Backend pobiera lub odczytuje dostępne dane dla lokalizacji i okresu. Raport uwzględnia także braki danych i ich pochodzenie.
  \item Moduł pogodowy wyznacza ewapotranspirację referencyjną \(\mathrm{ET}_0\) metodą FAO-56 Penmana--Monteitha, a w opisanych przypadkach braków promieniowania stosuje przybliżenie Hargreavesa. Zapotrzebowanie uprawy szacowane jest jako \(\mathrm{ETc}=K_c\times\mathrm{ET}_0\).
  \item Dobowy model strefy korzeniowej szacuje całkowity zapas wody dostępnej dla roślin (\(\mathrm{TAW}\)) i próg łatwo dostępnej wody (\(\mathrm{RAW}\)), uwzględniając retencję gleby, głębokość korzeni, opad i podaną dawkę nawadniania. Dwa warianty zapasu początkowego pokazują wrażliwość oszacowania.
  \item Wskaźniki Sentinel-2 (\code{NDVI}, \code{NDMI}, \code{NDWI}) dostarczają kontekstu o roślinności i wodzie powierzchniowej. Analizy rastrów stosują maskę jakości sceny i minimalny udział użytecznych pikseli.
  \item Dostępne dane terenowe i hydrologiczne uzupełniają ocenę. W paszporcie analiza NMT może opisywać zagłębienia, a przy dostępnych danych zlewni algorytm D8 wyznacza kierunek i akumulację spływu.
  \item Deterministyczne reguły backendu dobierają działania na podstawie bilansu, warunków działki i dostępności źródeł. Raport przedstawia podstawę rekomendacji; nie wyświetla gotowej, wcześniej przygotowanej oceny jako wyniku bieżącej analizy.
\end{enumerate}

\section{Architektura i technologie}

\begin{longtable}{@{}p{0.23\textwidth}p{0.71\textwidth}@{}}
\toprule
\rowcolor{PWForest}\textcolor{white}{\textbf{Obszar}} & \textcolor{white}{\textbf{Technologie i odpowiedzialność}} \\
\midrule
\endfirsthead
\toprule
\rowcolor{PWForest}\textcolor{white}{\textbf{Obszar}} & \textcolor{white}{\textbf{Technologie i odpowiedzialność}} \\
\midrule
\endhead
Interfejs & Angular 21 i TypeScript; PrimeNG 21, PrimeIcons i Tailwind CSS 4. Widoki obejmują analizę działki, paszport, eksperymenty i opis metod. \\
Mapa i wykresy & MapLibre GL JS, Terra Draw, Turf oraz Apache ECharts przez \code{ngx-echarts}. \\
API & Python 3.13, FastAPI, Pydantic i Uvicorn. \code{backend/main.py} udostępnia HTTP i WebSocket oraz może serwować zbudowany frontend. \\
Pobieranie i przetwarzanie & HTTPX, NumPy, Rasterio, Shapely i PyProj. Moduły pogodowe, przestrzenne i paszportu integrują źródła danych. \\
Analiza i raporty & \code{science.py} zawiera model pogodowy i glebowy; \code{analysis.py} składa ocenę i rekomendacje; \code{passport.py} buduje paszport; \code{field_analysis.py} analizuje gotowy paszport. ReportLab generuje PDF. \\
Indeks i przechowywanie & DuckDB obsługuje lokalny indeks działek. SQLite przechowuje wyniki pobrań, dane pochodne i raporty; raport po identyfikatorze jest dostępny do 14 dni. \\
Uruchomienie & Docker Compose uruchamia aplikację i osobną usługę \code{parcels-updater} do budowania i odświeżania indeksu działek. \\
Testy & Pytest dla backendu, Vitest dla frontendu i Playwright dla głównego przepływu w przeglądarce. \\
\bottomrule
\end{longtable}

W opisie projektu nie wskazano wykorzystania generatywnej sztucznej inteligencji. Opisane rekomendacje są wynikiem deterministycznych reguł i obliczeń modelowych.

\section{Źródła danych}

\begin{longtable}{@{}p{0.26\textwidth}p{0.68\textwidth}@{}}
\toprule
\rowcolor{PWForest}\textcolor{white}{\textbf{Źródło}} & \textcolor{white}{\textbf{Zakres wykorzystania}} \\
\midrule
\endfirsthead
\toprule
\rowcolor{PWForest}\textcolor{white}{\textbf{Źródło}} & \textcolor{white}{\textbf{Zakres wykorzystania}} \\
\midrule
\endhead
GUGiK / EGiB i krajowy GeoParquet & Geometrie i identyfikatory działek. Lokalny indeks DuckDB obejmuje według opisu projektu ok. 38,9 mln działek; \code{parcels-updater} buduje i odświeża indeks. ULDK może uzupełniać geometrię i dane administracyjne. \\
NASA POWER Daily AG & Dobowe szeregi opadu, temperatury, punktu rosy, wiatru i promieniowania wykorzystywane w bilansie pogodowym. \\
IMGW-PIB & Dane pobliskiej stacji do porównania warunków meteorologicznych; dla bieżącej analizy może być dostępny odczyt pobliskiego wodowskazu. \\
Copernicus Sentinel-2 L2A, Microsoft Planetary Computer i Earth Search & Katalogi STAC i fragmenty rastrów COG do analizy wskaźników satelitarnych. Główna analiza korzysta z obu katalogów, a paszport z Earth Search. \\
ISRIC SoilGrids 2.0 & WCS z modelowanymi właściwościami gleby w siatce 250 m, m.in. retencją przy 33 i 1500 kPa na kilku głębokościach. \\
Open Topo Data / EU-DEM, Copernicus DEM GLO-30 i GUGiK NMT & Wysokość, nachylenie, opis terenu i -- przy dostępności -- wykrywanie zagłębień oraz analiza spływu. \\
ESA WorldCover 2021 & Udziały klas pokrycia terenu i wykrywanie pikseli stałej wody w promieniu do 2 km od działki. \\
OpenStreetMap i Geoportal GUGiK & Kafelki mapy i ortofotomapa WMTS. W środowisku hackathonu nakładka WMS granic EGiB była niedostępna; wyszukiwanie działek korzysta z lokalnego indeksu. \\
Open-Meteo & Opcjonalne dane historyczne ERA5-Land, prognoza w paszporcie i awaryjne źródło pogody dla analizy. Ustawienie \code{WODO_NONCOMMERCIAL_OPEN_METEO=1} włącza opisany w projekcie bezpłatny dostęp niekomercyjny. \\
Copernicus Data Space / Sentinel Hub i Climate Data Store & Opcjonalna roczna klasyfikacja upraw CLMS z pewnością klasyfikacji oraz szeregi ERA5-Land wilgotności gleby na głębokości 100--289 cm. Integracje wymagają danych dostępowych; dla ERA5-Land przewidziano alternatywnie lokalny CSV. \\
Pliki operatora i dane rolnika & BDOT10k do odległości od cieków, rowów i stawów; NMT zlewni do analizy spływu; lokalna referencja Sentinel-2 do anomalii dla tej samej działki, uprawy i fazy; deklaracje rolnika jako dodatkowy kontekst. \\
\bottomrule
\end{longtable}

\section{Aktualność i pochodzenie danych}

Wynik i rekomendacja są obliczane przez aplikację na podstawie danych dla wskazanej działki i daty. Nie są gotowym zestawem wcześniej zapisanych rekomendacji. Dane i wyniki pobrań mogą być przechowywane lokalnie, a raport po identyfikatorze działki jest dostępny do 14 dni.

\begin{longtable}{@{}p{0.22\textwidth}p{0.30\textwidth}p{0.42\textwidth}@{}}
\toprule
\rowcolor{PWForest}\textcolor{white}{\textbf{Tryb}} & \textcolor{white}{\textbf{Dane}} & \textcolor{white}{\textbf{Informacja o dacie}} \\
\midrule
\endfirsthead
\toprule
\rowcolor{PWForest}\textcolor{white}{\textbf{Tryb}} & \textcolor{white}{\textbf{Dane}} & \textcolor{white}{\textbf{Informacja o dacie}} \\
\midrule
\endhead
Indeks lokalny przygotowywany i odświeżany przez usługę & Krajowy GeoParquet EGiB w DuckDB & Opis projektu nie podaje daty pobrania ani ostatniego odświeżenia indeksu. \\
Zewnętrzne usługi i API & NASA POWER, IMGW, katalogi Sentinel-2, SoilGrids, usługi NMT, OpenStreetMap / Geoportal oraz opcjonalnie Open-Meteo & Opis projektu nie określa harmonogramu pobierania ani reguł cache dla każdego źródła. Raport powinien prezentować datę analizy, pochodzenie i braki danych. \\
Produkt o określonej edycji & ESA WorldCover 2021 & Rok produktu jest znany; opis nie podaje daty jego pobrania do środowiska projektu. \\
Pliki lokalne lub deklaracje użytkownika & BDOT10k, NMT, referencja Sentinel-2 i dane gospodarstwa & Daty pozyskania plików i czas podania deklaracji nie zostały podane w materiale źródłowym. \\
Raporty i dane pochodne & SQLite & Opis podaje dostępność raportu do 14 dni, ale nie podaje szczegółowej polityki retencji pozostałych rekordów. \\
\bottomrule
\end{longtable}

Przed zgłoszeniem należy uzupełnić datę pozyskania lub aktualizacji lokalnego GeoParquet i plików operatora oraz potwierdzić, że raport zapisuje rzeczywisty czas pobrania danych ze źródeł zewnętrznych. Nie należy przedstawiać brakujących dat jako znanych.

\section{Interfejsy API}

FastAPI udostępnia dokumentację interaktywną pod \code{/docs} oraz specyfikację OpenAPI pod \code{/openapi.json}.

\begin{description}[style=nextline,leftmargin=0pt,labelwidth=0pt,labelsep=0pt,font=\normalfont\bfseries]
  \item[\code{GET /health}, \code{GET /api/health}, \code{GET /api/crops}] Stan aplikacji i indeksu, zakres dat i lista profili upraw.
  \item[\code{GET /api/parcels/{parcel_id}}, \code{GET /api/parcels/search}, \code{GET /api/parcels/by-point}] Wyszukiwanie działek po identyfikatorze, tekście lub punkcie.
  \item[\code{POST /api/parcels/import}] Przyjęcie geometrii GeoJSON.
  \item[\code{POST /api/analysis} i \code{WS /api/ws/analysis}] Uruchomienie analizy; WebSocket przekazuje postęp i wynik albo błąd.
  \item[\code{GET /api/analysis/{analysis_id}} i \code{GET /api/analysis/{analysis_id}/provenance}] Odczyt raportu oraz informacji o źródłach, brakach, dacie analizy i wersji modelu.
  \item[\code{GET /api/analysis/{analysis_id}/report?format=pdf}] Pobranie raportu w PDF lub JSON.
  \item[\code{GET /api/passport}, \code{POST /api/passport}, \code{POST /api/analysis/passport}] Utworzenie paszportu i analiza paszportu dla wybranej daty.
  \item[\code{GET /api/data-sources/status}] Status źródeł wykorzystanych w ostatniej analizie.
\end{description}

\section{Uruchomienie}

Opis projektu wskazuje Docker Compose jako sposób uruchomienia aplikacji i usługi \code{parcels-updater}, ale nie podaje polecenia startowego, nazw plików Compose, wymaganych zmiennych środowiskowych ani procedury inicjalizacji indeksu. Te szczegóły należy przepisać z konfiguracji wdrożeniowej właściwego repozytorium przed przekazaniem dokumentacji jury; nie można ich wiarygodnie odtworzyć z samego opisu funkcjonalnego.

Integracje z Sentinel Hub i Climate Data Store wymagają danych dostępowych. Open-Meteo w opisanym wariancie niekomercyjnym wymaga ustawienia \code{WODO_NONCOMMERCIAL_OPEN_METEO=1} i przestrzegania warunków dostawcy. Dostępność głównych funkcji w środowisku wdrożeniowym zależy od osiągalności zewnętrznych usług oraz skonfigurowania opcjonalnych źródeł.

\section{Dostępność dla rolnika}

Podstawowy przepływ odbywa się w aplikacji webowej: rolnik wybiera lub wskazuje działkę, podaje uprawę i datę, a aplikacja przygotowuje raport. Nie wymaga to od użytkownika znajomości GIS, programowania ani samodzielnego przetwarzania danych przestrzennych. Materiał źródłowy nie potwierdza testami zgodności z konkretnymi przeglądarkami ani nie określa kosztów utrzymania i ewentualnych odpłatnych warstw zewnętrznych; te informacje trzeba sprawdzić dla docelowego wdrożenia.

\section{Licencje i warunki korzystania}

Opis projektu wskazuje, że warunki wykorzystania i atrybucje zbiorów są dokumentowane osobno. W przekazanym materiale nie ma jednak treści tych wykazów ani pełnych warunków licencji bibliotek, danych i usług. Dlatego przed publikacją należy dołączyć i sprawdzić dla faktycznie użytych wersji: źródło, licencję lub regulamin, wymagane atrybucje, ograniczenia komercyjne oraz ewentualne limity API. Szczególnej weryfikacji wymaga wariant niekomercyjny Open-Meteo i opcjonalne usługi wymagające danych dostępowych.

\section{Ograniczenia i niepewność}

\begin{itemize}
  \item Dane pogodowe z siatki opisują obszar większy niż pojedyncze pole.
  \item Właściwości gleby SoilGrids i zapas wody w strefie korzeniowej są oszacowaniami, a warianty zapasu początkowego wpływają na wynik.
  \item Na obserwacje Sentinel-2 wpływają m.in. zachmurzenie, faza rozwoju i płodozmian; historia typu uprawy wymaga potwierdzenia przez użytkownika.
  \item Dla analiz historycznych współczesna geometria i statyczne produkty mogą być użyte jako oznaczony kontekst, nie jako historyczny pomiar.
  \item Model nie został jeszcze zwalidowany terenowymi pomiarami wilgotności z tych samych działek. Rekomendacje należy traktować jako wsparcie kontroli i planowania, nie jako gwarancję plonu ani samodzielną podstawę do określenia dawki nawodnienia.
  \item Ocena lokalnego odpływu i decyzja o robotach terenowych wymagają weryfikacji warunków w gospodarstwie. Ustalenie dawki nawadniania wymaga dodatkowych informacji o polu i dostępnych zasobach.
  \item Awaria, ograniczenia lub brak danych z usług zewnętrznych mogą ograniczyć zakres raportu; braki powinny być komunikowane zamiast zastępowane pozornie precyzyjnym wynikiem.
\end{itemize}

\section{Skalowalność i bezpieczeństwo}

Lokalny indeks obejmujący według opisu około 38,9 mln działek stanowi podstawę wyszukiwania w skali kraju. DuckDB obsługuje indeks, a SQLite cache i raporty. Przechowywanie wyników i dostępność raportu do 14 dni ograniczają część powtarzalnych obliczeń. Przy wdrożeniu dla wielu użytkowników należy dodatkowo monitorować limity i dostępność zewnętrznych usług, kontrolować liczbę zapytań, określić harmonogram aktualizacji danych i przetestować czas odpowiedzi oraz koszty infrastruktury. Materiał źródłowy nie podaje zmierzonych czasów odpowiedzi ani prognozy kosztów produkcyjnych.

Aplikacja nie powinna zbierać danych osobowych, jeśli nie są potrzebne do analizy. Deklaracje dotyczące gospodarstwa należy ograniczyć do informacji niezbędnych do kontekstu obliczeń. Opis projektu nie określa mechanizmów uwierzytelniania, zarządzania sekretami, kopii zapasowych ani szczegółowej polityki retencji; przed wdrożeniem publicznym te elementy wymagają osobnego potwierdzenia i konfiguracji.

\section{Testy i weryfikacja}

Opis projektu wymienia Pytest dla backendu, Vitest dla frontendu i Playwright dla głównego przepływu w przeglądarce. Nie podaje wyników uruchomienia tych testów, dlatego należy dołączyć aktualny raport z CI lub lokalnego wykonania.

Obliczenia mają deterministyczne reguły, a raport i endpoint pochodzenia udostępniają źródła, braki danych, datę analizy i wersję modelu. Walidacja naukowa na terenowych pomiarach wilgotności z tych samych działek nie została jeszcze przeprowadzona. Weryfikacja rekomendacji powinna więc obejmować przegląd zastosowanych równań i progów, kontrolę pochodzenia danych oraz porównanie z pomiarami polowymi, gdy będą dostępne.

\section{Przykład demonstracyjny}

Przykładowa analiza dotyczy działki \code{301103_2.0009.549/7} w gminie Kościan, uprawy ziemniaka i daty \code{2024-07-05}. Typ uprawy dla historycznego terminu wymaga potwierdzenia przez użytkownika.

\section{Informacje do uzupełnienia przed przekazaniem jury}

\begin{itemize}
  \item Dokładna instrukcja uruchomienia, pliki konfiguracyjne Compose, zmienne środowiskowe i sposób przygotowania indeksu.
  \item Daty pobrania i ostatniej aktualizacji lokalnych zbiorów oraz potwierdzenie rejestrowania dat pobrania danych z API.
  \item Pełny wykaz licencji, warunków usług i wymaganych atrybucji dla faktycznie użytych bibliotek, danych i komponentów.
  \item Wyniki testów i status wdrożenia prototypu, w tym działający adres lub instrukcja dostępu dla jury.
  \item Rzeczywisty raport wygenerowany przez prototyp oraz demonstracyjny film i prezentacja PDF, wymagane w zgłoszeniu.
\end{itemize}

\end{document}
